\chapter{Tailwind CSS for MERN Apps}

\begin{notebox}[NOTE]
This chapter assumes you already know CSS and focuses only on applying
Tailwind's utility classes to build the Task Management System's screens.
\end{notebox}

\section{Responsive Prefixes}

A prefix applies a style from a breakpoint up: \texttt{sm:} (640px),
\texttt{md:} (768px), \texttt{lg:} (1024px). Unprefixed classes apply at
all sizes.

\begin{lstlisting}
<div className="flex flex-col md:flex-row gap-4">
  {/* stacked on mobile, side-by-side from tablet up */}
</div>
\end{lstlisting}

\section{Flex, Grid, and Spacing}

\begin{lstlisting}
<div className="flex items-center justify-between p-4">
  <h1 className="text-xl font-bold">Tasks</h1>
  <button className="px-4 py-2 bg-blue-600 text-white rounded">New Task</button>
</div>

<div className="grid grid-cols-1 sm:grid-cols-2 lg:grid-cols-3 gap-4">
  {tasks.map((t) => <TaskCard key={t._id} task={t} />)}
</div>
\end{lstlisting}

Common spacing: \texttt{p-4}, \texttt{m-4}, \texttt{gap-4},
\texttt{space-y-4} (gap between stacked children with no flex/grid parent).

\section{Styling Forms}

\begin{lstlisting}
<label className="block text-sm font-medium text-gray-700 mb-1">Email</label>
<input type="email" className="w-full px-3 py-2 border border-gray-300 rounded-md
  focus:outline-none focus:ring-2 focus:ring-blue-500" />
\end{lstlisting}

\texttt{focus:ring-*} gives a visible focus ring for branding while
keeping keyboard accessibility.

\section{Buttons: Primary / Secondary / Danger}

\begin{lstlisting}[language=JavaScript]
// components/Button.jsx
const base = "px-4 py-2 rounded-md font-medium transition-colors disabled:opacity-50";
const variants = {
  primary: "bg-blue-600 text-white hover:bg-blue-700",
  secondary: "bg-gray-100 text-gray-800 hover:bg-gray-200",
  danger: "bg-red-600 text-white hover:bg-red-700",
};

function Button({ variant = "primary", children, ...props }) {
  return <button className={`${base} ${variants[variant]}`} {...props}>{children}</button>;
}
export default Button;
// <Button variant="danger">Delete</Button>
\end{lstlisting}

\section{Cards: TaskCard}

\begin{lstlisting}[language=JavaScript]
// components/TaskCard.jsx
function TaskCard({ task }) {
  const statusColor = {
    pending: "bg-yellow-100 text-yellow-800",
    "in-progress": "bg-blue-100 text-blue-800",
    done: "bg-green-100 text-green-800",
  };
  return (
    <div className="border rounded-lg p-4 shadow-sm hover:shadow-md bg-white">
      <div className="flex justify-between items-start mb-2">
        <h3 className="font-semibold">{task.title}</h3>
        <span className={`text-xs px-2 py-1 rounded-full ${statusColor[task.status]}`}>{task.status}</span>
      </div>
      <p className="text-sm text-gray-600 line-clamp-2">{task.description}</p>
    </div>
  );
}
export default TaskCard;
\end{lstlisting}

\section{Modals}

\begin{lstlisting}[language=JavaScript]
// components/Modal.jsx
function Modal({ isOpen, onClose, children }) {
  if (!isOpen) return null;
  return (
    <div className="fixed inset-0 bg-black/50 flex items-center justify-center z-50">
      <div className="bg-white rounded-lg shadow-xl w-full max-w-md p-6 relative">
        <button onClick={onClose} className="absolute top-3 right-3">&times;</button>
        {children}
      </div>
    </div>
  );
}
export default Modal;
// <Modal isOpen={showModal} onClose={() => setShowModal(false)}><TaskForm onSaved={...} /></Modal>
\end{lstlisting}

\begin{notebox}[NOTE]
\texttt{fixed inset-0} pins the overlay to the viewport; \texttt{bg-black/50}
is Tailwind's slash opacity syntax --- no separate opacity class needed.
\end{notebox}

\section{Navbar with Responsive Menu}

\begin{lstlisting}[language=JavaScript]
// components/Navbar.jsx
import { useState } from "react";

function Navbar() {
  const [open, setOpen] = useState(false);
  return (
    <nav className="bg-white border-b px-4 py-3">
      <div className="flex items-center justify-between">
        <span className="font-bold text-blue-600">TaskManager</span>
        <div className="hidden md:flex gap-6">
          <a href="/tasks">Tasks</a>
          <a href="/profile">Profile</a>
        </div>
        <button className="md:hidden" onClick={() => setOpen(!open)}>Menu</button>
      </div>
      {open && (
        <div className="md:hidden mt-3 flex flex-col gap-2">
          <a href="/tasks">Tasks</a>
          <a href="/profile">Profile</a>
        </div>
      )}
    </nav>
  );
}
export default Navbar;
\end{lstlisting}

\section{Sidebar Layout}

\begin{lstlisting}
<div className="flex min-h-screen">
  <aside className="hidden lg:block w-64 bg-gray-900 text-white p-4">{/* nav links */}</aside>
  <main className="flex-1 p-6 bg-gray-50">{/* page content */}</main>
</div>
\end{lstlisting}

\texttt{hidden lg:block} hides the sidebar below 1024px so mobile gets a
toolbar/nav instead of a permanent sidebar.

\begin{tipbox}[TIP]
\texttt{hidden md:flex} (or \texttt{lg:block}) plus a mobile toggle button
is the most reused responsive pattern across navbars, sidebars, and filters.
\end{tipbox}

\whatwhyhow
{Tailwind utility classes (layout, spacing, color, state variants, responsive prefixes) are composed directly in JSX to build buttons, cards, modals, navbars, and layouts.}
{Utility classes avoid separate CSS files and specificity fights; responsive prefixes adapt one component across breakpoints without media queries.}
{Build small components (\texttt{Button}, \texttt{TaskCard}, \texttt{Modal}) with variant classNames, and combine \texttt{hidden}/\texttt{flex}/\texttt{grid} with breakpoint prefixes for mobile-to-desktop layouts.}
